\section{Amplituden- und Frequenz- und Phasenumtastung (ASK, FSK, PSK)}
\label{section:ask_fsk_afsk}
\begin{frame}{Digitale Modulationsverfahren}
    \begin{itemize}
        \item Genauso wie es verschiedene analoge Modulationsverfahren gibt, gibt es auch verschiedene digitale Modulationsverfahren.
        \item Die grundlegenden Möglichkeiten ein Signal zu modulieren, also auf einen Hochfrequenzträger aufzuprägen, sind dieselben: Veränderung der Amplitude, der Frequenz oder der Phase des Trägers.
        \item Beim unmodulierten Träger hingegen bleiben Amplitude, Frequenz und Phasenlage konstant.
        \item Insbesondere bietet sich hier das I/Q-Verfahren an (dazu später mehr).
    \end{itemize}
\end{frame}

\begin{frame}{Amplitude Shift Keying, ASK}
    \begin{figure}
        \DARCimage{0.85\linewidth}{700include}
        \caption{\scriptsize Amplitudenumtastung (Amplitude-shift Keying)}
        \label{ask}
    \end{figure}
\end{frame}

\begin{frame}{Frequency Shift Keying, FSK}
    \begin{figure}
        \DARCimage{0.85\linewidth}{703include}
        \caption{\scriptsize Frequenzumtastung (Frequency-shift Keying)}
        \label{fsk}
    \end{figure}
\end{frame}

\begin{frame}{Phase Shift Keying, PSK}
    \begin{figure}
        \DARCimage{0.85\linewidth}{705include}
        \caption{\scriptsize Phasenumtastung (Phase-shift Keying)}
        \label{psk}
    \end{figure}
\end{frame}

\begin{frame}
\only<1>{
\begin{question2x2}{EE406}{Welches der folgenden Diagramme zeigt einen erkennbar durch Amplitudenumtastung (ASK) modulierten Träger?}{\DARCimage{1.0\linewidth}{357include}}
{\DARCimage{1.0\linewidth}{356include}}
{\DARCimage{1.0\linewidth}{358include}}
{\DARCimage{1.0\linewidth}{359include}}
\end{question2x2}

}
\only<2>{
\begin{question2x2}{EE406}{Welches der folgenden Diagramme zeigt einen erkennbar durch Amplitudenumtastung (ASK) modulierten Träger?}{\DARCimage{1.0\linewidth}{357include}}
{\DARCimage{1.0\linewidth}{356include}}
{\textbf{\textcolor{DARCgreen}{\DARCimage{1.0\linewidth}{358include}}}}
{\DARCimage{1.0\linewidth}{359include}}
\end{question2x2}

}
\end{frame}

\begin{frame}
\only<1>{
\begin{question2x2}{EE407}{Welches der folgenden Diagramme zeigt einen erkennbar durch Frequenzumtastung (FSK) modulierten Träger?}{\DARCimage{1.0\linewidth}{357include}}
{\DARCimage{1.0\linewidth}{356include}}
{\DARCimage{1.0\linewidth}{358include}}
{\DARCimage{1.0\linewidth}{359include}}
\end{question2x2}

}
\only<2>{
\begin{question2x2}{EE407}{Welches der folgenden Diagramme zeigt einen erkennbar durch Frequenzumtastung (FSK) modulierten Träger?}{\textbf{\textcolor{DARCgreen}{\DARCimage{1.0\linewidth}{357include}}}}
{\DARCimage{1.0\linewidth}{356include}}
{\DARCimage{1.0\linewidth}{358include}}
{\DARCimage{1.0\linewidth}{359include}}
\end{question2x2}

}
\end{frame}

\begin{frame}
\only<1>{
\begin{question2x2}{AE401}{Welches der folgenden Diagramme zeigt einen erkennbar durch Phasenumtastung (PSK) modulierten Träger?}{\DARCimage{1.0\linewidth}{357include}}
{\DARCimage{1.0\linewidth}{356include}}
{\DARCimage{1.0\linewidth}{359include}}
{\DARCimage{1.0\linewidth}{358include}}
\end{question2x2}

}
\only<2>{
\begin{question2x2}{AE401}{Welches der folgenden Diagramme zeigt einen erkennbar durch Phasenumtastung (PSK) modulierten Träger?}{\DARCimage{1.0\linewidth}{357include}}
{\DARCimage{1.0\linewidth}{356include}}
{\textbf{\textcolor{DARCgreen}{\DARCimage{1.0\linewidth}{359include}}}}
{\DARCimage{1.0\linewidth}{358include}}
\end{question2x2}

}
\end{frame}

\begin{frame}
\only<1>{
\begin{question2x2}{EE101}{Welches der folgenden Diagramme zeigt einen unmodulierten Träger?}{\DARCimage{1.0\linewidth}{356include}}
{\DARCimage{1.0\linewidth}{357include}}
{\DARCimage{1.0\linewidth}{359include}}
{\DARCimage{1.0\linewidth}{358include}}
\end{question2x2}

}
\only<2>{
\begin{question2x2}{EE101}{Welches der folgenden Diagramme zeigt einen unmodulierten Träger?}{\textbf{\textcolor{DARCgreen}{\DARCimage{1.0\linewidth}{356include}}}}
{\DARCimage{1.0\linewidth}{357include}}
{\DARCimage{1.0\linewidth}{359include}}
{\DARCimage{1.0\linewidth}{358include}}
\end{question2x2}

}
\end{frame}%ENDCONTENT
